\subsubsection{Uncertainty-based Continual Learning (UCL)}

\gls{ucl} \cite{ll_ucl_2019} is a Bayesian regularisation approach to \gls{cl} that replaces deterministic weight tensors with factorised Gaussian posteriors.

\paragraph{Bayesian Parameterisation}

Every trainable parameter in the network (convolutional and linear) is represented by two real-valued parameter tensors: a mean $\mu$ and a raw standard-deviation encoding $\rho$. The posterior standard deviation is recovered via the softplus transform,
%
\begin{equation}\label{eqn:ucl_sigma}
    \sigma = \log\!\left(1 + e^{\rho}\right),
\end{equation}
%
ensuring $\sigma > 0$ for all $\rho \in \mathbb{R}$. During a stochastic forward pass a weight sample is drawn using the reparameterisation trick,
%
\begin{equation}\label{eqn:ucl_reparam}
    w = \mu + \sigma \cdot \varepsilon,
    \qquad
    \varepsilon \sim \mathcal{N}(0, I),
\end{equation}
%
while the deterministic path uses $w = \mu$ directly.

\paragraph{Initialisation}

Both convolutional and linear Bayesian layers are initialised to approximate Kaiming He initialisation for the mean while targeting a prescribed noise level for the variance. Given layer fan-in $n$, a noise ratio $r \in (0,1)$ is applied to split the total He variance $v_\text{total} = 2/n$ into
%
\begin{equation}\label{eqn:ucl_init}
    v_\mu = v_\text{total}(1-r),
    \qquad
    v_\epsilon = v_\text{total} \cdot r,
\end{equation}
%
so that $\mu$ is drawn uniformly from $[-\sqrt{3 v_\mu},\, \sqrt{3 v_\mu}]$ and $\rho$ is initialised to $\log(\exp(\sqrt{v_\varepsilon}) - 1)$, i.e.\ the softplus inverse of $\sqrt{v_\varepsilon}$.

% \paragraph{Architecture}

% The model consists of three components.

% \subparagraph{Bayesian ResNet-18 backbone.} A \texttt{BayesianResNet1D} replaces every \texttt{Conv1d} in the standard ResNet-18 stack with a \texttt{BayesianConv1d} that maintains $(\mu, \rho)$ tensors instead of a single weight tensor. Batch normalisation layers are kept deterministic; their running statistics are updated during training. The network produces a $512$-dimensional feature vector via global average pooling.

% \subparagraph{Per-task Bayesian heads.} A \texttt{ModuleList} of $T$ independent \texttt{BayesianLinear} layers maps the feature vector to per-task class logits. Each head $k$ has output dimension $c_k$ (the number of classes for task $k$).

% \subparagraph{Task-incremental vs.\ class-incremental operation.} When \texttt{split=True} (task-incremental, \gls{til}), the forward pass returns only the logits from head $t$. When \texttt{split=False} (class-incremental, \gls{cil}), all heads are concatenated along the class dimension before returning, yielding a single logit vector of width $\sum_k c_k$.

\paragraph{Optimisation}

Mean parameters $\{\mu\}$ and adapter parameters are optimised at learning rate $\eta_\mu$; variance parameters $\{\rho\}$ at a separate, typically larger, rate $\eta_\rho$. Both groups are passed to a single \texttt{Adam} optimiser with two parameter groups. During each \texttt{observe} call, the model executes $I_\text{inner}$ gradient steps (default $I_\text{inner}=1$), each drawing a fresh weight sample via \eqref{eqn:ucl_reparam}.

\paragraph{UCL Regularisation}

At task-transition time, a frozen snapshot of the current posterior $(\mu^\text{old}, \sigma^\text{old})$ is stored as a \texttt{BayesianClassifier} clone with gradients disabled. During subsequent tasks, a regularisation penalty $\mathcal{R}$ is added to the classification loss, accumulating contributions from the backbone layers and from all previous task heads.

For each matched layer pair (old snapshot vs.\ current trainable layer), define the per-parameter importance as
%
\begin{equation}\label{eqn:ucl_strength}
    s_i = \frac{\sigma^\text{init}_i}{\sigma^\text{old}_i},
    \qquad
    \sigma^\text{init}_i = \sqrt{\frac{2}{n} \cdot r},
\end{equation}
%
where $n$ is the fan-in of the layer and $r$ is the noise ratio. Layers whose posterior has sharpened since initialisation ($\sigma^\text{old}_i < \sigma^\text{init}_i$) therefore impose a stronger penalty on weight drift.

Six raw penalty terms are accumulated per layer:
%
\begin{align}
    R_{\sigma} &= \sum_i
        \left(
            \frac{(\sigma^\text{new}_i)^2}{(\sigma^\text{old}_i)^2}
            - \log\!\left(\frac{(\sigma^\text{new}_i)^2}{(\sigma^\text{old}_i)^2}\right)
        \right),
        \label{eqn:ucl_sigma_reg} \\
    R_{\sigma,0} &= \sum_i
        \left(
            (\sigma^\text{new}_i)^2
            - \log\!\left((\sigma^\text{new}_i)^2 \right)
        \right),
        \label{eqn:ucl_sigma_normal_reg} \\
    R_{\mu,W} &= \sum_i
        \left(s_i \left(\mu^\text{new}_i - \mu^\text{old}_i\right)\right)^2,
        \label{eqn:ucl_mu_weight_reg} \\
    R_{\mu,b} &= \sum_j
        \left(\bar{s}_j \left(b^\text{new}_j - b^\text{old}_j\right)\right)^2,
        \label{eqn:ucl_mu_bias_reg} \\
    R_{L1,W} &= \sum_i
        \frac{(\mu^\text{old}_i)^2}{(\sigma^\text{old}_i)^2}
        \left|\mu^\text{new}_i - \mu^\text{old}_i\right|,
        \label{eqn:ucl_l1_weight_reg} \\
    R_{L1,b} &= \sum_j
        \frac{(b^\text{old}_j)^2}{\bar\sigma^2_j}
        \left|b^\text{new}_j - b^\text{old}_j\right|,
        \label{eqn:ucl_l1_bias_reg}
\end{align}
%
where $\bar{s}_j$ and $\bar\sigma_j$ are the row-mean importance and standard deviation for row $j$, obtained by averaging $s_i$ and $\sigma^\text{old}_i$ over the input dimension. All six terms are normalised by the total number of regularised parameters $N_\text{reg}$ before being combined.

The final per-minibatch loss is
%
\begin{equation}\label{eqn:ucl_loss}
    \mathcal{L}_\text{UCL}
    =
    \mathcal{L}_\text{CE}
    + \frac{\alpha}{2B}\!\left(R_{\mu,W} + R_{\mu,b}\right)
    + \frac{1}{B}\!\left(R_{L1,W} + R_{L1,b}\right)
    + \frac{\beta}{2B}\!\left(R_{\sigma} + R_{\sigma,0}\right),
\end{equation}
%
where $B$ is the minibatch size, $\alpha$ controls mean-drift penalties, and $\beta$ controls variance-drift penalties. The L1 terms are scaled by the binary flag \texttt{saved} (1 after the first task transition, 0 during the first task), so they activate only once a snapshot exists.

\paragraph{Interpretation}
\begin{figure}
    \centering
    \includegraphics[width=0.75\linewidth]{03_Related Work//Figures/ucl_log_var_penalty.png}
    \caption{UCL variance-drift term versus variance ratio $r=\sigma_{\text{new}}^{2}/\sigma_{\text{old}}^{2}$, showing the full penalty $r-\log(r+\epsilon)$ (solid) and log component $-\log(r+\epsilon)$ (dashed). The dashed vertical line at $r=1$ marks unchanged variance; arrows indicate variance decrease ($r<1$) and increase ($r>1$).}
    \label{fig:ucl_log-var-penalty}
\end{figure}

The variance penalty $R_\sigma$ in \eqref{eqn:ucl_sigma_reg} penalises ratios $(\sigma^\text{new})^2 / (\sigma^\text{old})^2$ that deviate from 1, with the log correction making the penalty asymmetric: it is cheaper to sharpen the posterior than to broaden it. Intuitively, this constrains the network from ``unlearning'' high-confidence weights on old tasks. The normal penalty $R_{\sigma,0}$ in \eqref{eqn:ucl_sigma_normal_reg} independently penalises each current variance $(\sigma^\text{new}_i)^2$ for deviating from 1---it anchors the variance from changing too sporadically and/or towards extreme values.

The quadratic mean penalty $R_{\mu,W}$ in \eqref{eqn:ucl_mu_weight_reg} is an importance-weighted analogue of the \gls{ewc} quadratic term \eqref{eqn:ewc_loss}: here the importance $s_i = \sigma^\text{init}/\sigma^\text{old}$ is derived from the ratio of initial to current posterior uncertainty rather than from squared gradients. Weights whose posterior has concentrated (small $\sigma^\text{old}$) are penalised more strongly for drifting. The L1 penalty $R_{L1,W}$ in \eqref{eqn:ucl_l1_weight_reg} provides an additional, sparse penalty proportional to the old-posterior signal-to-noise ratio $\mu^2/\sigma^2$, i.e.\ parameter magnitude normalised by uncertainty, so it captures confidence in a nonzero effect size (high for large-and-certain weights, low for small and/or uncertain weights) that variance alone cannot represent.

\paragraph{Inference}

At test time, the model can either use the posterior mean ($\mu$) directly, or average predictions over $S$ stochastic weight samples (Monte Carlo integration). When $S > 1$, each sample draws weights via \eqref{eqn:ucl_reparam}, computes class probabilities via softmax, and the mean predictive distribution
%
\begin{equation}\label{eqn:ucl_mc}
    \bar{p}(y \mid x) = \frac{1}{S} \sum_{s=1}^{S} p(y \mid x, w^{(s)})
\end{equation}
%
is used to form log-probability logits. This enables epistemic uncertainty estimation via mutual information between the predictive distribution and the weight posterior. In \gls{til} mode the relevant task head is selected by task identity $t$; in \gls{cil} mode, a logit mask is applied to zero out classes beyond those seen up to the current task, and the full concatenated vector is used for prediction.
